\begin{figure}[ht]
\centering
\begin{tikzpicture}[x=1cm,y=1cm,>=stealth]
\def\cw{0.55}   % cell width
\def\rh{1.05}   % row height
% a row of three cells at height #1, lit pattern #2#3#4 (1 = lit), ringed pivot column #5 (0 = none)
\newcommand{\cellrow}[5]{%
  \foreach \c/\lit in {1/#2,2/#3,3/#4}{%
    \ifnum\lit=1 \fill[black!75] ({(\c-1)*\cw},#1) rectangle ({\c*\cw},{#1+\cw}); \fi
    \draw ({(\c-1)*\cw},#1) rectangle ({\c*\cw},{#1+\cw});
  }%
  \ifnum#5>0 \draw[line width=1pt] ({(#5-0.5)*\cw},{#1+0.5*\cw}) circle ({0.38*\cw}); \fi
}
% column headers: the points 1, 2, 3 of ray 1
\foreach \c in {1,2,3}{\node[font=\scriptsize] at ({(\c-0.5)*\cw},{5*\rh+\cw+0.25}) {$\c$};}
% the two lamps
\cellrow{5*\rh}{1}{0}{1}{0}
\node[anchor=east,font=\small] at (-0.3,{5*\rh+0.5*\cw}) {$u=\{1,3\}$};
\node[anchor=west,font=\small] at ({3*\cw+0.3},{5*\rh+0.5*\cw}) {$T_u=\tau_1\tau_3$};
\cellrow{4*\rh}{0}{1}{0}{0}
\node[anchor=east,font=\small] at (-0.3,{4*\rh+0.5*\cw}) {$w=\{2\}$};
\node[anchor=west,font=\small] at ({3*\cw+0.3},{4*\rh+0.5*\cw}) {$T_w=\tau_2$};
% step 1: one translation word
\draw[->] ({1.5*\cw},{4*\rh-0.1}) -- ({1.5*\cw},{3*\rh+\cw+0.1});
\node[anchor=west,font=\footnotesize] at ({3*\cw+0.3},{3.5*\rh+0.5*\cw}) {Step 1: $T_u^{-1}T_w\to T_{u_0}$; two relations \textsf{Q}, two swaps};
\cellrow{3*\rh}{1}{1}{1}{3}
\node[anchor=east,font=\small] at (-0.3,{3*\rh+0.5*\cw}) {$u_0=u+w=\{1,2,3\}$};
\node[anchor=west,font=\small] at ({3*\cw+0.3},{3*\rh+0.5*\cw}) {$T_{u_0}=\tau_1\tau_2\tau_3$, pivot $3$};
% step 2, first transvection: from the pivot of row u_0 into cell 1 of row u_1
\draw[->,thick] ({2.5*\cw},{3*\rh-0.05}) .. controls ({2.5*\cw},{2*\rh+0.5*\cw}) and ({0.5*\cw},{3*\rh-0.1}) .. ({0.5*\cw},{2*\rh+\cw+0.05});
\node[anchor=west,font=\footnotesize] at ({3*\cw+0.3},{2.5*\rh+0.5*\cw}) {Step 2: $X_{13}$ adds $3$ into $1$; one swap, one relation \textsf{Q}};
\cellrow{2*\rh}{0}{1}{1}{3}
\node[anchor=east,font=\small] at (-0.3,{2*\rh+0.5*\cw}) {$u_1=\{2,3\}$};
\node[anchor=west,font=\small] at ({3*\cw+0.3},{2*\rh+0.5*\cw}) {$T_{u_1}=\tau_2\tau_3$};
% step 2, second transvection: from the pivot of row u_1 into cell 2 of row u_2
\draw[->,thick] ({2.5*\cw},{2*\rh-0.05}) .. controls ({2.5*\cw},{1*\rh+0.5*\cw}) and ({1.5*\cw},{2*\rh-0.1}) .. ({1.5*\cw},{1*\rh+\cw+0.05});
\node[anchor=west,font=\footnotesize] at ({3*\cw+0.3},{1.5*\rh+0.5*\cw}) {$X_{23}$ adds $3$ into $2$; one relation \textsf{Q}};
\cellrow{1*\rh}{0}{0}{1}{3}
\node[anchor=east,font=\small] at (-0.3,{1*\rh+0.5*\cw}) {$u_2=e_3$};
\node[anchor=west,font=\small] at ({3*\cw+0.3},{1*\rh+0.5*\cw}) {$T_{u_2}=\tau_3$};
% step 3
\node[anchor=west,font=\footnotesize] at ({3*\cw+0.3},{0.5*\rh+0.5*\cw-0.15}) {Step 3: $G=X_{23}X_{13}$ is linear; $[x,\tau_3y\tau_3^{-1}]$ is a local filling};
\end{tikzpicture}
\caption{The clearing of the worked example, $N=3$: the lit cells are the support of a configuration on the points $1,2,3$ of ray~$1$, the ringed
cell is the pivot, and each arrow is a transvection that adds the value at the pivot into the cell it points to. After $X_{13}$ and $X_{23}$ the
configuration $u_0=\{1,2,3\}$ has become $e_3$, and the translation word has been renormalized after each step.}
\label{fig:clearing}
\end{figure}
